\documentclass[11pt]{article}
\usepackage{lmodern}          % AER font
\usepackage[T1]{fontenc}      % Better font encoding
\usepackage{amsmath, amssymb}
\usepackage[margin=1in]{geometry}
\usepackage{graphicx}
\usepackage{url}
\usepackage[round]{natbib}
\usepackage{bibentry}         % lets \nobibliography resolve cites without printing a list
\usepackage{hyperref}

%=======================================
% Format helpers
%=======================================
% Target school -- used in the header and the conclusion so they stay in sync.
\newcommand{\school}{[University name]}
\newcommand{\program}{[Program name]}

% Run-in section heading: bold label ending in a period, body text follows
% on the same line (matches the reference format).
\newcommand{\shead}[1]{\par\vspace{0.85em}\noindent\textbf{#1.}\ }

% Indented paragraphs, no extra vertical space between them.
\setlength{\parindent}{1.5em}
\setlength{\parskip}{0pt}


\begin{document}

%=======================================
% Header
%=======================================
\noindent
\begin{minipage}[t]{0.60\textwidth}
  {\large\bfseries Statement of Purpose}
\end{minipage}\hfill
\begin{minipage}[t]{0.38\textwidth}
  \raggedleft
  Rhea Choudhary\\[2pt]
  \school
\end{minipage}

\vspace{0.5em}
\noindent\rule{\textwidth}{0.8pt}

\vspace{1.1em}

%=======================================
% Body
%=======================================
\shead{Research interests} I am applying to pursue a PhD in \program at \school, where I intend to study questions at the intersection of public, labour and health economics. My specific research agenda asks how large-scale government programs can better achieve their desired policy objectives, and the externalities created by such programs in the pursuit of these goals. \\~\\
So far, I have begun to tackle this research agenda in three key ways. Firstly, in my undergraduate Honours thesis published as a central bank working paper, I evaluated the effects of unconventional monetary policy -- which constitutes large central bank investment in instilling market confidence -- on financial stability in an emerging market context. Secondly, through my pre-doctoral research I am examining how the introduction of Australia's National Disability Insurance Scheme (NDIS) changed the labour market composition and returns to disability work, and potential effects on care. Finally, in my independent research, I have been exploring the efficacy of large-scale supply-side education policies on improving individual employment outcomes as well as labour market shortages. 
Methodologically, I seek to use microdata, particularly Australian microdata which I have used as part of my pre-doctoral research, with theoretical models to imbue greater intuition to empirical findings.

\shead{Pre-doctoral and professional experience} I have been broadening my research toolkit working for Professor Manasi Deshpande as a pre-doctoral researcher, where we have been exploiting the staggered rollout of NDIS to examine how this large-scale government program funding disability services causally changed labour supply in disability work, wages, related professions such as nursing and therapy, and categorise the returns to NDIS among workers, firms and households. Through this research I have gained significant experience working with large administrative datasets, applying an array of causal inference methods, as well as implementing a LASSO model to identify a clean `risk set' of individuals for whom NDIS may have induced labour supply. During my pre-doctoral position, I audited Professor Deshpande's PhD elective course on Empirical Topics in Social Insurance, where I learned about the challenges in policy design, targeting and heterogeneity in outcomes, but also about innovative techniques to generate variation, assign treatment, and study these programs in ways that meaningfully contribute to the literature. Overall, my pre-doctoral experience has taught me how to think deeply, creatively and methodically about complex economics questions, and to conduct economic research deliberately and credibly. Prior to my pre-doctoral position, I worked at the Reserve Bank of Australia (RBA), Australia's equivalent of the Federal Reserve. In this role, I led an array of analytical research, most recently applying contingent valuation and textual analysis in a large-scale population representative survey, which is in the process of becoming a central bank working paper.

\shead{Future research} While auditing Professor Deshpande's course on the empirical analysis of social insurance, I became interested in how social transfers target cohorts when need is unobservable \citep{akerlof1978,nichols1982}, and the challenge in balancing efficient targeting with creating barriers to take-up among intended cohorts \citep{deshpande_li2019}. I was also curious about how I could apply learnings from this literature in social insurance to government programs more broadly. An Australian education policy called `Free TAFE', where from 2019-2022 the state of Victoria offered a free vocational education for certain courses that were deemed to have labour shortages, could be an interesting case-study to study this topic.\footnote{Examples include disability care, nursing, construction, and early childhood care. From 2023 onwards, this program was phased in by other states and received Commonwealth funding.} This policy has the unique characteristic that eligibility is governed by what courses are available (at the `good' level) rather than at the recipient level (e.g., by means-testing or employment history). 
\\~\\
Therefore, in this research, I ask two questions: firstly, whether targeting on the subsidised good produces different outcomes from participant-side targeting, and more generally, whether preferentially subsidised training is a legitimate lever for upskilling individuals into growth sectors when implemented at scale. I have been using linked unit-level administrative data from the Australian Bureau of Statistics (ABS), which allow me to observe individuals' prior education history, socio-economic background, parental income, health, and employment outcomes over time.\\~\\
To answer my first question on targeting, I leverage the fact that free TAFE was an unanticipated program and one initially offered only in the state of Victoria. I compare whether the program changed the composition of enrolments in these courses across age, gender, labour force status at enrolment, and highest education level. Inspired by the targeting efficiency work in \citet{deshpande_li2019} and complier analysis from \citet{mountjoy2022}, I intend to explore how well this course-based targeting measure fared in inducing likely `target cohorts' to enter the program (e.g., low-skilled, low-income cohorts), dissuaded individuals with low marginal returns to this education from enrolling, and how much of the enrolment effect comes from inframarginal enrolees.
\\~\\
On my second research question, I am interested in how labour supply and wages adjusted to this policy in general equilibrium. For example, I am interested in whether wage premia in targeted sectors changed if supply of skilled workers increased, and whether provision shifts the composition of the local training market. In my work thus-far, I find that low-skilled women have the strongest enrolment first-stage, and so as extensions I would like to consider how outcomes like marriage and fertility changed in response to this program.
%=======================================

% Citations resolve against bib.bib, but no reference list is printed.
\bibliographystyle{abbrvnat}
\nobibliography{bib}

\end{document}
